% TOP_PROBLEM: {"id": 30006738, "title": "Regular Inverse Galois Problem over Q", "category_id": 4, "category": {"id": 4, "name": "algebra", "display_name": "Algebra", "description": "Group theory, ring theory, field theory, and algebraic structures.", "slug": "algebra", "order_index": 4, "created_at": "2026-07-31T15:26:25.671Z"}, "status": "open"}
% FIRST_POSTED: 2026-09-27
% !TeX program = lualatex
% Compile twice: lualatex 30006738.tex
% All references and prose are embedded; no BibTeX database or external inputs.
\documentclass[11pt,a4paper]{article}
\usepackage[margin=24mm,headheight=14pt]{geometry}
\usepackage{fontspec}
\setmainfont{Latin Modern Roman}
\setsansfont{Latin Modern Sans}
\setmonofont{Latin Modern Mono}
\usepackage{amsmath,amssymb,mathtools}
\usepackage{unicode-math}
\setmathfont{Latin Modern Math}
\usepackage{microtype}
\usepackage{enumitem,xurl,needspace}
\usepackage[dvipsnames]{xcolor}
\usepackage{hyperref}
\hypersetup{unicode=true,colorlinks=true,linkcolor=MidnightBlue,urlcolor=MidnightBlue,
 pdftitle={500 Mathematical Problem Statements},pdfauthor={Source-annotated compilation},bookmarksnumbered=true}
\usepackage{bookmark}
\usepackage{tocloft}
\setlength{\cftsecnumwidth}{3em}
\cftsetpnumwidth{2.5em}
\cftsetrmarg{4em}
\usepackage{titlesec}
\titleformat{\section}{\Large\bfseries\raggedright}{\thesection}{0.7em}{}
\usepackage{fancyhdr}
\pagestyle{fancy}\fancyhf{}
\fancyhead[L]{\small\sffamily Mathematical problem statements}
\fancyhead[R]{\small Source edition 22}
\fancyfoot[C]{\thepage}
\setlength{\parindent}{0pt}
\setlength{\parskip}{5pt plus 1pt}
\setlength{\emergencystretch}{3em}
\setcounter{tocdepth}{1}
\setcounter{secnumdepth}{1}
\setlist[itemize]{leftmargin=3em,itemsep=5pt,topsep=4pt,parsep=0pt}
\allowdisplaybreaks
\newcommand{\N}{\mathbb{N}}
\newcommand{\Z}{\mathbb{Z}}
\newcommand{\Q}{\mathbb{Q}}
\newcommand{\R}{\mathbb{R}}
\newcommand{\C}{\mathbb{C}}
\newcommand{\F}{\mathbb{F}}
\newcommand{\A}{\mathbb{A}}
\newcommand{\PP}{\mathbb P}
\newcommand{\E}{\mathbb{E}}
\newcommand{\eps}{\varepsilon}
\newcommand{\dd}{\,\mathrm d}
\newcommand{\sref}[2]{\hyperlink{source:#1:#2}{[S#2]}}
\newcommand{\eref}[2]{\hyperlink{extra:#1:#2}{[E#2]}}
% Turn the next switch off for a shorter reading copy. The quotations preserve
% every exactTarget field, including qualifications omitted from a short summary.
\newif\ifshowcatalogue
\showcataloguetrue
\newcommand{\cataloguescope}[1]{\ifshowcatalogue
 \paragraph{Exact scope in the supplied catalogue.}
 \begin{quote}\small #1\end{quote}\fi}
\hypersetup{pdftitle={Problem 30006738 - Regular inverse Galois problem},pdfauthor={Open Problems Math research notebook}}
\fancyhead[L]{\small\sffamily Open Problems Math \textbar\ Research notebook}
\fancyhead[R]{\small\sffamily 120 \textbar\ 27 September 2026}
\numberwithin{equation}{section}
\begin{document}
\setcounter{section}{0}
% ================================================================
% problemId: 30006738
% Canonical problem ID: 30006738
% exactTarget SHA-256: 094c9c84eaa9404e559586387f22b5b43d93ebf7f5bfe04d71bc2787323c98df
\clearpage
\section*{\texorpdfstring{Regular Inverse Galois Problem over Q}{Regular Inverse Galois Problem over Q}}\label{30006738}
\subsection{Definitions and mathematical statement}
A finite extension $F/\Q(T)$ is regular over $\Q$ when $F\cap\overline{\Q}=\Q$ inside an algebraic closure, so there are no new algebraic constants. The question is
\[
 \forall G\text{ finite}\ \exists F/\Q(T)\text{ finite Galois and regular}:\quad
 \operatorname{Gal}(F/\Q(T))\simeq G.
\]
Geometrically, this asks for a geometrically integral Galois cover of $\PP^1_{\Q}$ with group $G$. A realization only over $\Q$ or with a nontrivial constant-field extension does not meet the regular target.
\par\smallskip\textit{Formulation and concept references:} \sref{30006738}{1}.
\cataloguescope{For every finite group G, does there exist a finite Galois extension F/Q(T) with Gal(F/Q(T)) isomorphic to G and Q algebraically closed in F? Equivalently, require a geometrically integral Galois cover of the projective line over Q with group G. A positive resolution proves this for all finite groups; a negative resolution proves nonexistence for at least one finite group.}
\subsection{Short English statement}
Can every finite group be realized by a regular Galois cover of the projective line defined over the rationals?
% BEGIN RESEARCH ADDITION 120
\subsection{Research attempt}

\paragraph{Current frontier.}
The formulation source supplies new regular realizations for certain
$\operatorname{PSL}_2(\F_{p^2})$ groups subject to arithmetic conditions;
it also distinguishes regular realizations from realizations only over
$\Q$.  Its introduction and main results, checked in the December 2025
version, do not give all finite groups. \sref{30006738}{1}
The constructions below are independent elementary deductions illustrating
what an all-groups argument must preserve: geometric connectedness and
no new constants.  A complex cover, a number-field extension, or a cover
of a higher-genus base is not silently substituted.

\paragraph{Attack routes.}
A branch-cycle approach can start with topological covers over $\C$, but
needs a rational moduli point and valid descent to $\Q$.  A generic
permutation action can produce every finite group over an invariant
function field, but must reduce that base to $\Q(T)$.  A constructive
closure approach combines already realizable groups.  We prove the
closure operations carefully, then examine the first operation that
fails to yield a universal construction.

\paragraph{Attempt I: a fully checked regular example.}
Consider the splitting field $F$ of
\[
 f(X)=X^3-X-T\in\Q(T)[X].
\]
Over any characteristic-zero field $k$, the inclusion
$k(T)\subset k(x)$ with $T=x^3-x$ has degree three: the rational map
$\PP^1_x\to\PP^1_T$ has degree three, or equivalently the pole order
of $T$ at infinity is three.  Thus $f$ is irreducible over both $\Q(T)$
and $\overline\Q(T)$.  Its discriminant is
\[
 \operatorname{disc}_X(f)=4-27T^2.
\]
This has two simple zeros over $\overline\Q$, so is not a square in
$\overline\Q(T)$.  An irreducible cubic has transitive Galois group in
$S_3$; if that group were $A_3$, its alternating product of root
differences would be fixed, making the discriminant a square.  Conversely
that product detects the sign character.  Hence the geometric and
arithmetic groups are both $S_3$.

Let $k_0=F\cap\overline\Q$.  The geometric group is the kernel of the
arithmetic action on $k_0$ and here already has order six, the order of
the full arithmetic group.  Thus $k_0=\Q$.  This proves a regular
$S_3$ realization with the catalogue's exact constant-field requirement.
The discriminant identity is also verified symbolically by the
companion script; no finite specialization is used to infer regularity.

\paragraph{Attempt II: closure under quotients and products.}
If $F/\Q(T)$ is regular Galois with group $G$ and $N\triangleleft G$,
then $F^N/\Q(T)$ is Galois with group $G/N$, and
$F^N\cap\overline\Q\subseteq F\cap\overline\Q=\Q$.
Thus regular realizability is closed under quotients.

It is also closed under finite direct products.  To see this, take two
regular Galois covers with groups $G_1,G_2$.  Their finite branch sets
in $\PP^1(\overline\Q)$ can be made disjoint by applying a rational
M\"obius transformation to one base.  Indeed, for each pair of branch
points the condition that the transformation send one to the other is
a proper algebraic condition on $\operatorname{PGL}_2$.  Finitely many
such conditions do not exhaust its rational points: on a rational affine
chart their product is a nonzero polynomial, and $\Q$ is infinite.
This also handles branch points at infinity by using suitable charts.

After base extension to $\overline\Q(T)$, the intersection of the two
Galois function fields is Galois and can ramify only in the intersection
of their branch sets.  It is therefore unramified everywhere.  The
Riemann--Hurwitz formula for a connected unramified degree-$d$ cover of
$\PP^1$ gives $2g-2=-2d$, impossible for $d>1$.  The intersection is
trivial.  The two fields are consequently linearly disjoint and their
geometric compositum has group $G_1\times G_2$.

The arithmetic compositum embeds into the same product, whereas its
geometric subgroup already has the full product order.  Equality
follows, with constant field $\Q$.  Iteration proves finite-product
closure.  For example, every quotient of a finite product of copies
of the displayed $S_3$ extension has a regular realization.  This is
an actual family, but not a family known to contain every finite group.

\paragraph{Attempt III: why the evident subgroup construction stops.}
Every finite group embeds in a symmetric group.  Suppose we have a
regular $S_m$ extension $F/\Q(T)$ and a subgroup $H\leq S_m$.
The extension $F/F^H$ has group $H$, but its base is $F^H$, not generally
$\Q(T)$.  Its corresponding curve need not have genus zero.  Regularity
of $F/\Q$ ensures that the base has no extra algebraic constants; it
says nothing about that curve being rational.  The operation therefore
does not prove closure under arbitrary subgroups in the required
one-variable sense.  The preceding quotient argument works precisely
because the base stays fixed.

A universal version of the same issue begins with independent variables
$\{x_g:g\in G\}$.  Let $G$ act by left translation on their labels.
The action on $L=\Q(x_g:g\in G)$ is faithful, and the elementary fixed
field theorem for finite automorphism groups yields
\[
 L/L^G\text{ Galois with group }G.
\]
Both fields have constant field $\Q$: an element of a purely
transcendental field algebraic over $\Q$ belongs to $\Q$.  However
$\operatorname{trdeg}_{\Q}L^G=|G|$, so this is not a realization over
$\Q(T)$.  Setting many variables to rational functions of one parameter
can identify branches or reduce the Galois group.

Here is an exact sufficient geometric target.  Restrict the linear
permutation representation to its open set $U$ of points with trivial
stabilizer.  The quotient map $U\to Y=U/G$ is a finite \emph{\'{e}tale}
$G$-torsor in characteristic zero.  It would suffice to find, for each
$G$, a nonempty open subset $V\subset\PP^1_{\Q}$ and a morphism
$V\to Y$ whose pulled-back torsor is geometrically connected.  The
function field of that pullback is then regular Galois over $\Q(T)$
with group $G$: connectedness after base change makes the geometric
monodromy transitive on a regular $G$-fiber, hence equal to $G$.
This is a sufficient rational-curve-and-monodromy condition, not a claim
that every regular realization must arise from this particular model.

Rational points of $Y$ alone are insufficient.  Their fibers give
finite extensions over $\Q$, not one geometrically connected cover
varying over $V$.  Likewise a rational curve in $Y$ whose image lies in
a locus of smaller monodromy does not suffice.  These distinctions are
exactly where a naive specialization argument runs out of justification.

\paragraph{Stress test.}
The product argument used disjoint branch sets \emph{geometrically},
not merely disjoint sets of rational branch points.  It used
characteristic zero to exclude nontrivial everywhere-unramified covers
of the geometric projective line; it did not claim an analogous argument
for arbitrary higher-genus bases.  The $S_3$ example checks geometric
irreducibility as well as the arithmetic discriminant, preventing an
unnoticed constant extension.

The missing rational-curve condition is substantially stronger than
having many complex covers or a rational point on some moduli space.
No descent obstruction has been proved to vanish for every $G$, and no
group has been shown to have an unavoidable obstruction for every possible
cover.  In particular, failure of this chosen invariant-field strategy
would not disprove the original conjecture.  A positive solution would
also imply the ordinary inverse Galois problem over $\Q$ through Hilbert
irreducibility, as the formulation source explains; this consequence
underscores the depth of the unproved step. \sref{30006738}{1}

\paragraph{Outcome.}
\textbf{Outcome: Exploratory approach with unresolved gap.}

The notebook proves a regular $S_3$ example and the quotient/product
closure operations, then isolates a sufficient geometrically connected
pullback condition for a universal construction.  The direct subgroup
and generic-invariant-field shortcuts fail to preserve the required
base.  No all-groups construction or counterexample has been obtained;
the elementary lemmas are not claimed to be new to the literature.
% END RESEARCH ADDITION 120
\subsection{Sources}
\begin{itemize}\raggedright
\item[\textnormal{[S1]}]\hypertarget{source:30006738:1}{} Demeio and Gvirtz-Chen, Galois Realisations of PSL2(Fp\textasciicircum{}2).
\url{https://arxiv.org/html/2512.23615v1}.
{\small\textit{Catalogue formulation source.}}
\end{itemize}


\end{document}
